\documentclass[11pt,a4paper]{article}

% ── Packages ──
\usepackage[utf8]{inputenc}
\usepackage[T1]{fontenc}
\usepackage{amsmath,amssymb}
\usepackage{physics}
\usepackage{siunitx}
\usepackage{booktabs}
\usepackage{geometry}
\usepackage{hyperref}
\usepackage{enumitem}
\usepackage{float}

\geometry{margin=2.5cm}
\hypersetup{colorlinks=true,linkcolor=blue,citecolor=blue,urlcolor=blue}

\title{Residual Gas Velocity Above a Stopped Propeller Blade\\
       in Dense Xenon: A Mechanism-by-Mechanism Analysis}
\author{}
\date{}

\begin{document}
\maketitle

% ══════════════════════════════════════════════════════════════
\section{Problem Definition}
\label{sec:problem}
% ══════════════════════════════════════════════════════════════

A horizontal two-bladed propeller with NACA\,0012 airfoil cross-section
rotates inside a vertical cylindrical vessel filled with xenon gas at
high pressure. The propeller sweeps through a prescribed angle
$\theta_\mathrm{sweep}$ in a time $T_\mathrm{sweep}$, starting from
rest and returning to rest. We seek the residual gas velocity at a
height $z_\mathrm{obs}$ above the upper blade surface, evaluated at a
time $t_\mathrm{obs}$ after the blade has stopped.

\subsection{Parameters}

Seven parameters are treated as variable inputs; all others are fixed.
The rotation time $t_\mathrm{rot}$ is computed from the inertia and
motor torque (see Section~\ref{sec:kinematics}).

\begin{table}[H]
\centering
\caption{Variable parameters (with default values).}
\label{tab:params}
\begin{tabular}{llll}
\toprule
Symbol & Description & Default & Units \\
\midrule
$c$                   & Blade chord length          & 0.16  & m   \\
$\theta_\mathrm{sweep}$ & Sweep angle              & $\pi$ & rad \\
$m_\mathrm{blade}$   & Mass of each blade ($\times 2$)  & 0.25  & kg  \\
$m_\mathrm{plate}$   & Mass of each carrier plate ($\times 2$) & 0.15 & kg \\
$\tau$               & Motor torque                & 60    & N$\cdot$m \\
$z_\mathrm{obs}$     & Observation height above blade top & 0.005 & m \\
$t_\mathrm{obs}$     & Observation time after stop  & 0.5   & s   \\
\bottomrule
\end{tabular}
\end{table}

\begin{table}[H]
\centering
\caption{Fixed parameters.}
\label{tab:fixed}
\begin{tabular}{lll}
\toprule
Parameter & Value & Description \\
\midrule
$R$            & \SI{1.6}{m}        & Blade span (hub to tip) \\
$P$            & \SI{15}{bar}       & Xenon pressure \\
$T$            & \SI{300}{K}        & Gas temperature \\
NACA profile   & 0012               & Symmetric, 12\% thickness \\
$\alpha$       & $0^\circ$          & Angle of attack \\
$d_\mathrm{wire}$  & \SI{0.2}{mm}  & Wire grid wire diameter \\
$s_\mathrm{wire}$  & \SI{2}{mm}    & Wire grid spacing (centre-to-centre) \\
$\Delta D$     & \SI{0.10}{m}       & Vessel inner diameter surplus over propeller \\
\bottomrule
\end{tabular}
\end{table}

\subsection{Geometry and Coordinate System}

The blade rotates in the $xy$-plane. The $z$-axis points upward from
the blade centre-plane. The NACA\,0012 blade has maximum thickness
$t_\mathrm{max} = 0.12\,c$, so the upper blade surface is at
$z_\mathrm{blade} = t_\mathrm{max}/2 = 0.06\,c$ above the rotation
axis. A horizontal wire grid sits at height $z_\mathrm{grid}$ above
the blade top (in this analysis we evaluate conditions at an arbitrary
$z_\mathrm{obs}$). The cylindrical vessel has inner radius
\begin{equation}
  R_\mathrm{cyl} = R + \frac{\Delta D}{2}\,.
\end{equation}
The tip clearance is $\Delta D / 2$.


% ══════════════════════════════════════════════════════════════
\section{Xenon Gas Properties}
\label{sec:gas}
% ══════════════════════════════════════════════════════════════

Xenon is a monatomic ideal gas to good approximation at \SI{15}{bar},
\SI{300}{K} (compressibility factor $Z \approx 0.98$).

\paragraph{Density.}
\begin{equation}
  \rho = \frac{P M_\mathrm{Xe}}{R_g T}\,,
  \label{eq:density}
\end{equation}
where $M_\mathrm{Xe} = \SI{0.13129}{kg/mol}$ and
$R_g = \SI{8.314}{J/(mol\cdot K)}$. This gives
$\rho \approx \SI{79}{kg/m^3}$, roughly 65 times denser than air at STP.

\paragraph{Dynamic viscosity.}
The viscosity of Xe is nearly pressure-independent at moderate densities.
At $T = \SI{300}{K}$:
\begin{equation}
  \mu \approx \SI{2.32e-5}{Pa\cdot s}\,.
\end{equation}

\paragraph{Kinematic viscosity.}
\begin{equation}
  \nu = \frac{\mu}{\rho} \approx \SI{2.94e-7}{m^2/s}\,.
  \label{eq:nu}
\end{equation}
This is extremely small --- about 50 times smaller than air at STP ---
because $\rho$ is large while $\mu$ is comparable to air.

\paragraph{Speed of sound.}
For a monatomic ideal gas with $\gamma = 5/3$:
\begin{equation}
  c_s = \sqrt{\frac{\gamma R_g T}{M_\mathrm{Xe}}}
      \approx \SI{178}{m/s}\,.
  \label{eq:cs}
\end{equation}


% ══════════════════════════════════════════════════════════════
\section{Blade Kinematics and Flow Regime}
\label{sec:kinematics}
% ══════════════════════════════════════════════════════════════

\subsection{Bang-Bang Motion Profile with Aerodynamic Drag}

The propeller executes a \emph{bang-bang} motion profile: it accelerates
under constant motor torque $\tau$, then decelerates under reversed
torque $-\tau$. The motion is governed by:
\begin{equation}
  I\,\frac{d\omega}{dt} = \pm\tau - \tau_\mathrm{drag}(\omega)\,,
  \label{eq:EOM}
\end{equation}
where the upper sign applies during acceleration and the lower during
deceleration.

\paragraph{Moment of Inertia.}
The system consists of two blades (uniform rods of mass $m_\mathrm{blade}$
extending from $r = 0$ to $r = R$) and two carrier plates (point masses
$m_\mathrm{plate}$ at the worst-case position $r = R$):
\begin{equation}
  I = 2 \times \frac{1}{3}\,m_\mathrm{blade}\,R^2
    + 2 \times m_\mathrm{plate}\,R^2
    = \frac{2R^2}{3}\left(m_\mathrm{blade} + 3\,m_\mathrm{plate}\right).
  \label{eq:I}
\end{equation}
For default parameters ($m_\mathrm{blade} = \SI{0.25}{kg}$,
$m_\mathrm{plate} = \SI{0.15}{kg}$, $R = \SI{1.6}{m}$):
$I \approx \SI{1.19}{kg\cdot m^2}$.

\paragraph{Aerodynamic Drag Torque.}
The drag on both blades at angular velocity $\omega$ is:
\begin{equation}
  \tau_\mathrm{drag} = \mathcal{D}\,\omega^2\,,
  \qquad
  \mathcal{D} = \frac{\rho\, C_d\, d_\mathrm{wake}\, R^4}{4}\,,
  \label{eq:tau_drag}
\end{equation}
where $d_\mathrm{wake} = 0.12\,c$ is the wake thickness (approximately
equal to the blade thickness for a streamlined section), and
$C_d \approx 0.1$ is an effective drag coefficient for the
blade assembly moving through the gas. Note that the NACA\,0012 profile
drag coefficient at $\alpha = 0$ is much smaller
($C_d^\mathrm{profile} \approx 0.007$), but this value applies to flow
along the chord. The higher effective $C_d$ used here accounts for:
(i) the blade cross-section sweeping tangentially through the gas, and
(ii) the carrier plates ($m_\mathrm{plate}$) sliding on the blades,
which add semi-bluff frontal area and flow disruption at interfaces.
The combined blade-plus-plate assembly is thus intermediate between
a streamlined airfoil and a bluff body.

For default parameters ($\rho = \SI{79}{kg/m^3}$, $C_d = 0.1$,
$d_\mathrm{wake} = 0.12 \times \SI{0.16}{m} = \SI{19.2}{mm}$,
$R = \SI{1.6}{m}$):
\begin{equation}
  \mathcal{D} \approx \SI{0.25}{N\cdot m\cdot s^2/rad^2}\,.
\end{equation}

\paragraph{Asymmetric Switching Point.}
Since drag always opposes motion, deceleration (where both reversed
motor torque \emph{and} drag oppose motion) is stronger than
acceleration. To complete the full sweep angle $\theta$, the switching
point $\theta_\mathrm{switch}$ must satisfy $\theta_\mathrm{switch} > \theta/2$.

The equation of motion~\eqref{eq:EOM} is nonlinear and is solved
numerically to find $\theta_\mathrm{switch}$ such that
$\omega \to 0$ exactly when $\theta \to \theta_\mathrm{target}$.

\paragraph{Inertia-Only Reference (No Drag).}
For comparison, if drag is neglected:
\begin{align}
  \alpha_0 &= \frac{\tau}{I} \approx \SI{50}{rad/s^2}\,,
  \label{eq:alpha} \\
  t_\mathrm{rot,0} &= 2\sqrt{\frac{\theta\, I}{\tau}}
    \approx \SI{500}{ms}\,,
  \label{eq:trot0} \\
  \omega_\mathrm{max,0} &= \sqrt{\frac{\theta\,\tau}{I}}
    \approx \SI{12.6}{rad/s}\,.
  \label{eq:omega_max0}
\end{align}

\paragraph{Numerical Solution with Drag.}
For default parameters, the numerical integration yields:
\begin{align}
  \theta_\mathrm{switch} &\approx \SI{2.05}{rad}
    \quad (117.5^\circ, \text{ vs.\ } 90^\circ \text{ without drag})\,,
  \label{eq:theta_switch} \\
  t_\mathrm{accel} &\approx \SI{306}{ms}\,,
  \label{eq:taccel} \\
  t_\mathrm{decel} &\approx \SI{201}{ms}\,,
  \label{eq:tdecel} \\
  t_\mathrm{rot} &\approx \SI{507}{ms}
    \quad (+1.4\% \text{ compared to no-drag})\,,
  \label{eq:trot} \\
  \omega_\mathrm{max} &\approx \SI{11.8}{rad/s}
    \quad (-6.3\% \text{ reduction})\,.
  \label{eq:omega_max}
\end{align}

\paragraph{Peak Tip Velocity.}
\begin{equation}
  V_\mathrm{tip} = \omega_\mathrm{max}\,R
    \approx \SI{18.8}{m/s}\,.
  \label{eq:Vtip}
\end{equation}

\paragraph{Drag Impact Summary.}
Aerodynamic drag has a modest effect on the kinematics:
\begin{itemize}[nosep]
  \item Rotation time increases by $\sim 1.4\%$.
  \item Peak angular velocity decreases by $\sim 6.3\%$.
  \item The switching angle shifts from $90^\circ$ to $117.5^\circ$.
\end{itemize}
The qualitative behaviour remains unchanged, and the peak velocity is
still $\sim\SI{19}{m/s}$, well within the incompressible regime.

\subsection{Flow Regime}

We evaluate the flow regime at the peak velocity, which represents the
most demanding conditions for the boundary layer.

\paragraph{Mach Number.}
The tip Mach number is
\begin{equation}
  \mathrm{Ma}_\mathrm{tip}
    = \frac{V_\mathrm{tip}}{c_s}\,.
  \label{eq:Mach}
\end{equation}
For default parameters, $\mathrm{Ma} \approx 0.11 < 0.3$, so the flow is
\emph{incompressible}: density variations are negligible ($< 5\%$),
the continuity equation simplifies to $\nabla \cdot \mathbf{v} = 0$,
and no shock waves or compressibility-driven phenomena occur.

\paragraph{Reynolds Number.}
The chord-based Reynolds number at the tip is
\begin{equation}
  \mathrm{Re}_\mathrm{tip} = \frac{V_\mathrm{tip}\,c}{\nu}\,.
  \label{eq:Re}
\end{equation}
For default parameters, $\mathrm{Re}_\mathrm{tip} \approx 1.1 \times 10^7$,
far exceeding the critical value of $\sim 5 \times 10^5$ for transition.
The boundary layer is therefore \emph{fully turbulent}: inertial forces
dominate viscous forces, producing chaotic eddies with enhanced mixing
near the wall. The turbulent BL is thicker than a laminar one but has
a fuller velocity profile (more momentum near the wall).


% ══════════════════════════════════════════════════════════════
\section{Boundary Layer on the Upper Surface}
\label{sec:BL}
% ══════════════════════════════════════════════════════════════

Since we are interested in the flow \emph{above} the blade, only the
upper-surface boundary layer (BL) matters.

\subsection{Flat-Plate Turbulent BL Estimates}

For a turbulent BL developing from the leading edge over a flat plate
of length $c$ at Reynolds number $\mathrm{Re}_c = V c / \nu$, the
standard $\tfrac{1}{7}$-power-law result gives:
\begin{align}
  \delta &= 0.37\,c\;\mathrm{Re}_c^{-1/5}\,,
    \label{eq:delta} \\
  \theta &= \frac{7}{72}\,\delta
          = \frac{7}{72}\times 0.37\,c\;\mathrm{Re}_c^{-1/5}\,,
    \label{eq:theta} \\
  \delta^* &= \frac{1}{8}\,\delta\,,
    \label{eq:dstar}
\end{align}
where $\delta$ is the BL thickness, $\theta$ the momentum thickness,
and $\delta^*$ the displacement thickness.

\subsection{NACA\,0012 Pressure Gradient Effects}

An explicit integral BL calculation (Thwaites' method for the laminar
portion, Head's entrainment method for the turbulent portion) using the
NACA\,0012 inviscid velocity distribution shows that the favourable
pressure gradient over the front $\sim 30\%$ of the chord thins the BL,
while the adverse gradient over the rear $\sim 70\%$ thickens it. These
two effects approximately cancel: at $x/c \approx 0.95$ (the last
station before the integral method approaches separation),
$\delta \approx \SI{4.0}{mm}$, in excellent agreement with the
flat-plate estimate of $\SI{3.9}{mm}$. We therefore use the flat-plate
formula~\eqref{eq:delta} throughout.

\subsection{Skin Friction and Drag Coefficient}

The local skin-friction coefficient from the Blasius correlation:
\begin{equation}
  C_f = 0.0592\;\mathrm{Re}_c^{-1/5}\,.
  \label{eq:Cf}
\end{equation}

The friction velocity:
\begin{equation}
  u_\tau = V_\mathrm{tip}\,\sqrt{\frac{C_f}{2}}\,.
  \label{eq:utau}
\end{equation}

The profile drag coefficient for NACA\,0012 at $\alpha = 0$ is obtained
from the Squire--Young formula applied at the trailing edge, or from
experimental data:
\begin{equation}
  C_d \approx 0.006\text{--}0.009
  \label{eq:Cd}
\end{equation}
depending on the transition location. For the numerics we compute $C_d$
from the BL quantities via
\begin{equation}
  C_d = \frac{2\,(2\theta)}{c}
      \left(\frac{U_e(x_\mathrm{TE})}{V}\right)^{(H+5)/2},
  \label{eq:SY}
\end{equation}
where $H = \delta^*/\theta$ is the shape factor at the trailing edge.

\subsection{Wake Geometry (Upper Surface Only)}

The relevant length scales for default parameters are:
\begin{equation}
  \delta \approx \SI{2.3}{mm}\,,
  \quad
  \theta \approx \SI{0.23}{mm}\,,
  \quad
  \delta^* \approx \SI{0.29}{mm}\,.
\end{equation}
The upper blade surface sits at $z_\mathrm{blade} = 0.06\,c$ above the
rotation axis. The BL extends to
\begin{equation}
  z_\mathrm{BL} = 0.06\,c + \delta\,.
\end{equation}
The gap between the top of the BL and the observation height is
\begin{equation}
  \Delta z_\mathrm{gap} = z_\mathrm{obs} - \delta\,.
  \label{eq:gap}
\end{equation}
For default parameters ($z_\mathrm{obs} = \SI{5}{mm}$,
$\delta \approx \SI{2.3}{mm}$), $\Delta z_\mathrm{gap} \approx \SI{2.7}{mm}$.


\subsection{Initial Velocity Profile in the Lab Frame}

In the blade's reference frame, the BL velocity follows a
$\tfrac{1}{7}$-power law: $u_\mathrm{blade}(z) = V(z/\delta)^{1/7}$
for $0 \le z \le \delta$, with $u = V$ outside the BL. Transforming to
the lab frame (where the gas is initially at rest and the blade moves
at $V$):
\begin{equation}
  u_\mathrm{lab}(z) =
  \begin{cases}
    V\bigl[1 - (z/\delta)^{1/7}\bigr] & 0 < z \le \delta \\
    0 & z > \delta
  \end{cases}
  \label{eq:ulab}
\end{equation}
The gas near the blade surface ($z \to 0$) is dragged along at nearly
$V$; at $z = \delta$ the velocity is zero.


% ══════════════════════════════════════════════════════════════
\section{Mechanism 1: Turbulent Eddy Diffusion}
\label{sec:eddies}
% ══════════════════════════════════════════════════════════════

\subsection{Statement of the Problem}

After the blade stops, the momentum in the BL
(equation~\ref{eq:ulab}) can in principle diffuse vertically into the
quiescent gas above. The question is whether this diffusion carries any
appreciable velocity to $z = z_\mathrm{obs}$ within time
$t_\mathrm{obs}$.

The evolution is governed by the one-dimensional diffusion equation:
\begin{equation}
  \frac{\partial u}{\partial t}
  = \frac{\partial}{\partial z}
    \left[\nu_\mathrm{eff}(z,t)\,\frac{\partial u}{\partial z}\right],
  \label{eq:diffusion}
\end{equation}
with initial condition~\eqref{eq:ulab} and boundary conditions
\begin{equation}
  u(0,t) = 0 \quad\text{(no-slip on stationary blade)}\,,
  \qquad
  u(z \to \infty, t) = 0\,.
\end{equation}

\subsection{Turbulent Viscosity Profile}

The effective viscosity is $\nu_\mathrm{eff} = \nu + \nu_t(z,t)$.
At $t = 0$ (blade just stopped), the turbulent viscosity in the BL
follows the standard two-layer model:
\begin{equation}
  \nu_t(z, 0) =
  \begin{cases}
    \kappa\, u_\tau\, z & z < 0.2\,\delta
      \quad\text{(inner / log-law region)} \\
    0.018\, V\, \delta^* & 0.2\,\delta \le z < \delta
      \quad\text{(outer region)} \\
    0 & z \ge \delta
  \end{cases}
  \label{eq:nut_profile}
\end{equation}
where $\kappa = 0.41$ is the von K\'arm\'an constant. For default
parameters:
\begin{align}
  \nu_t^\mathrm{inner}(z = 0.2\delta)
    &= \kappa\, u_\tau\,(0.2\delta)
    \approx \SI{1.3e-4}{m^2/s}\,,
    \label{eq:nut_inner}
  \\
  \nu_t^\mathrm{outer}
    &= 0.018\, V\, \delta^*
    \approx \SI{1.0e-4}{m^2/s}\,.
    \label{eq:nut_outer}
\end{align}
The ratio to molecular viscosity is
$\nu_t / \nu \approx 350$.

\subsection{Turbulence Decay After Blade Stops}

Once the blade stops, there is no energy input to sustain the
turbulence. The eddies decay on the large-eddy turnover timescale:
\begin{equation}
  t_e = \frac{\delta}{u'}\,,
  \label{eq:teddy}
\end{equation}
where $u' \approx 0.05\,V$ is the outer-layer turbulence intensity
(comparable to $u_\tau$). For default parameters,
$t_e \approx \SI{2.3}{ms}$.

We model the decay as exponential with an $e$-folding time of
$2\,t_e$ (approximately two eddy turnovers for significant energy
loss):
\begin{equation}
  \nu_t(z,t) = \nu_t(z,0)\;\exp\!\left(-\frac{t}{2\,t_e}\right).
  \label{eq:nut_decay}
\end{equation}

\subsection{Diffusion Length Scale}

In the brief window before the turbulence dies ($\sim 2\,t_e \approx
\SI{5}{ms}$), the turbulent diffusion extends momentum by a distance
\begin{equation}
  \ell_\mathrm{turb} \sim \sqrt{\nu_t^\mathrm{outer} \times 2\,t_e}
  \approx \sqrt{\SI{1.0e-4}{} \times \SI{0.005}{}}
  \approx \SI{0.7}{mm}\,.
  \label{eq:ell_turb}
\end{equation}
Starting from the top of the BL at $z = \delta \approx \SI{2.3}{mm}$,
the momentum can reach at most $z \approx \SI{3.0}{mm}$. For
$z_\mathrm{obs} = \SI{5}{mm}$, the gap of \SI{2.7}{mm} is larger
than $\ell_\mathrm{turb}$.

After the turbulence decays, only molecular diffusion remains. In the
remaining time ($\sim \SI{0.5}{s}$):
\begin{equation}
  \ell_\mathrm{mol} = \sqrt{\nu\, t_\mathrm{obs}}
  \approx \sqrt{\SI{2.94e-7}{} \times \SI{0.5}{}}
  \approx \SI{0.38}{mm}\,.
  \label{eq:ell_mol}
\end{equation}
The total reach is $\delta + \ell_\mathrm{turb} + \ell_\mathrm{mol}
\approx \SI{3.4}{mm}$, short of $z_\mathrm{obs} = \SI{5}{mm}$.

\subsection{No-Slip Absorption}

Simultaneously, the no-slip condition at $z = 0$ (the now-stationary
blade surface) absorbs momentum from below. Since most of the momentum
is concentrated near $z = 0$ (equation~\ref{eq:ulab}), this absorption
is efficient. The numerical solution of equation~\eqref{eq:diffusion}
shows that $\sim 38\%$ of the initial momentum is absorbed by the
surface within \SI{0.5}{s}.

\subsection{Result}

The numerical solution of the 1D diffusion
equation~\eqref{eq:diffusion} with decaying turbulent viscosity gives:
\begin{equation}
  \boxed{u(z_\mathrm{obs},\, t_\mathrm{obs})
    \approx 0 \quad \text{(eddy diffusion)}}
\end{equation}
The turbulent eddies are too small ($\sim\delta = \SI{2.3}{mm}$) and
decay too quickly ($\sim\SI{5}{ms}$) to transport momentum across
the gap to $z_\mathrm{obs}$.


% ══════════════════════════════════════════════════════════════
\section{Mechanism 2: Bulk Swirl and Angular Momentum}
\label{sec:swirl}
% ══════════════════════════════════════════════════════════════

\subsection{Angular Impulse}

The drag on both blades deposits tangential (azimuthal) momentum into
the gas. The drag per unit span at radius $r$ is
\begin{equation}
  f(r) = \tfrac{1}{2}\,\rho\,[\omega r]^2\,c\,C_d\,.
\end{equation}
The total torque:
\begin{equation}
  \tau = 2 \int_0^R r\, f(r)\, dr
      = \rho\,\omega^2\,c\,C_d\,\frac{R^4}{4}\,.
  \label{eq:torque}
\end{equation}
The angular impulse over the sweep:
\begin{equation}
  J = \tau\, T_\mathrm{sweep}\,.
  \label{eq:J}
\end{equation}

\subsection{Where the Momentum Resides}

This angular momentum is deposited \emph{exclusively} within the
boundary layers on the blade surfaces --- a slab of thickness
$\sim(t_\mathrm{max} + 2\delta)$ centred on the blade plane. The gas
outside this slab is undisturbed (apart from the potential flow, which
vanishes when the blade stops; see Section~\ref{sec:potential}).

The key question is whether this momentum can propagate vertically to
$z_\mathrm{obs}$.

\subsection{Molecular Diffusion Timescale}

The time for momentum to diffuse across the gap
$\Delta z_\mathrm{gap}$ by molecular viscosity alone:
\begin{equation}
  t_\mathrm{visc}
    = \frac{(\Delta z_\mathrm{gap})^2}{\nu}
    = \frac{(z_\mathrm{obs} - \delta)^2}{\nu}\,.
  \label{eq:tvisc}
\end{equation}
For default parameters,
$t_\mathrm{visc} = (0.0027)^2 / \SI{2.94e-7}{} \approx \SI{25}{s}$
--- much longer than $t_\mathrm{obs} = \SI{0.5}{s}$.

\subsection{Ekman Pumping}

In a rotating fluid adjacent to a boundary, the Ekman layer drives a
secondary (radial then axial) circulation. The Ekman layer thickness
and the resulting axial velocity are:
\begin{equation}
  \delta_E = \sqrt{\frac{\nu_\mathrm{eff}}{\Omega}}\,,
  \qquad
  w_E = \sqrt{\nu_\mathrm{eff}\,\Omega}\,,
  \label{eq:Ekman}
\end{equation}
where $\Omega \sim u_\mathrm{wake}/R$ is the effective rotation rate of
the wake gas.

\paragraph{With molecular viscosity:}
$\delta_E \approx \SI{0.22}{mm}$, $w_E \approx \SI{1.4}{mm/s}$.
Negligible.

\paragraph{With turbulent viscosity} ($\nu_t \approx \SI{1e-4}{m^2/s}$):
$\delta_E \approx \SI{4}{mm}$, $w_E \approx \SI{2.6}{cm/s}$. But the
turbulence decays in $\sim\SI{5}{ms}$, so the vertical distance
traversed is $w_E \times 2 t_e \approx \SI{0.12}{mm}$.

In both cases, Ekman pumping cannot reach $z_\mathrm{obs}$ from the
blade plane.

\subsection{Ekman Spin-Up Time}

Even if one posits a mechanism for establishing vessel-scale secondary
circulation, the Ekman spin-up timescale is:
\begin{equation}
  t_\mathrm{Ekman} = \frac{H}{\sqrt{\nu_\mathrm{eff}\,\Omega}}\,,
  \label{eq:tEkman}
\end{equation}
where $H$ is the vessel height. With turbulent $\nu$:
$t_\mathrm{Ekman} \approx \SI{39}{s}$. With molecular $\nu$:
$t_\mathrm{Ekman} \approx \SI{740}{s}$. Both are far longer than
$t_\mathrm{obs}$.

\subsection{Result}

\begin{equation}
  \boxed{u(z_\mathrm{obs},\, t_\mathrm{obs})
    \approx 0 \quad \text{(bulk swirl)}}
\end{equation}
The angular momentum stays confined to the blade plane. No mechanism
transports it to $z_\mathrm{obs}$ within $t_\mathrm{obs}$.


% ══════════════════════════════════════════════════════════════
\section{Mechanism 3: Displacement (Potential) Flow}
\label{sec:potential}
% ══════════════════════════════════════════════════════════════

\subsection{Flow Field During the Sweep}

The NACA\,0012 blade has finite thickness
$t_\mathrm{max} = 0.12\,c$. As it moves through the gas, it displaces
fluid. The inviscid (potential) flow field around a two-dimensional body
of effective diameter $d \approx t_\mathrm{max}$ moving at velocity $V$
produces a vertical velocity component at height $z$ above the blade
centre that scales as:
\begin{equation}
  w_\mathrm{disp}(z)
    \sim V \,\frac{(d/2)^2}{z^2 + (d/2)^2}\,.
  \label{eq:wdisp}
\end{equation}
For $z = z_\mathrm{obs}$ and default parameters, this gives
$w_\mathrm{disp} \sim \SI{16}{m/s}$ --- a large value, present
\emph{only while the blade is in motion and in the vicinity}.

\subsection{Acoustic Clearing Time}

The potential flow field is slaved to the blade. When the blade stops,
the disturbance propagates away at the speed of sound:
\begin{equation}
  t_\mathrm{acoustic} = \frac{R}{c_s} \approx \SI{9}{ms}\,.
  \label{eq:tacoustic}
\end{equation}
Within $\sim\SI{10}{ms}$ of the blade stopping, the potential flow
field has been radiated away as a weak acoustic wave.

\subsection{Result}

\begin{equation}
  \boxed{u(z_\mathrm{obs},\, t_\mathrm{obs})
    = 0 \quad \text{(potential flow, for } t_\mathrm{obs}
    \gg t_\mathrm{acoustic}\text{)}}
\end{equation}


% ══════════════════════════════════════════════════════════════
\section{Mechanism 4: Pressure Pulse from Deceleration}
\label{sec:pressure}
% ══════════════════════════════════════════════════════════════

When the blade decelerates to rest over some time $t_\mathrm{decel}$,
the change in momentum of the fluid near the blade creates a pressure
transient.

\subsection{Amplitude}

The pressure perturbation scales as:
\begin{equation}
  \Delta p \sim \rho\, V_\mathrm{tip}^2\,.
  \label{eq:dp}
\end{equation}
For default parameters, $\Delta p \approx \SI{32000}{Pa}$, which is
$2\%$ of the ambient \SI{15}{bar} --- a small perturbation.

The velocity induced by this acoustic wave:
\begin{equation}
  \Delta u \sim \frac{\Delta p}{\rho\, c_s}
            = \frac{V_\mathrm{tip}^2}{c_s}\,.
  \label{eq:du_acoustic}
\end{equation}
For default parameters, $\Delta u \approx \SI{2.3}{m/s}$.

\subsection{Propagation}

The pressure pulse arrives at $z_\mathrm{obs}$ in a time
\begin{equation}
  t_\mathrm{arrive} = \frac{z_\mathrm{obs}}{c_s}\,.
\end{equation}
For default parameters, $t_\mathrm{arrive} \approx \SI{28}{\mu s}$.
The pulse has duration $\sim t_\mathrm{decel}$ and then propagates past
$z_\mathrm{obs}$, leaving no residual velocity.

\subsection{Result}

\begin{equation}
  \boxed{u(z_\mathrm{obs},\, t_\mathrm{obs})
    = 0 \quad \text{(pressure pulse, for }
    t_\mathrm{obs} \gg t_\mathrm{decel}\text{)}}
\end{equation}


% ══════════════════════════════════════════════════════════════
\section{Mechanism 5: Secondary Flows}
\label{sec:secondary}
% ══════════════════════════════════════════════════════════════

\subsection{Tip Clearance Flow}

In turbomachinery, a pressure difference between the suction and
pressure sides of a blade drives leakage flow through the tip gap. For
a lifting blade this is a significant source of axial mixing.

For NACA\,0012 at $\alpha = 0$: there is \emph{no lift}, and hence
\emph{no pressure difference} across the blade. The flow field is
symmetric about the chord plane.
\begin{equation}
  \Delta p_\mathrm{tip} = 0 \qquad (\alpha = 0)\,.
\end{equation}
No tip clearance flow exists.

\subsection{Centrifugal Effects and Wake--Wall Interaction}

The tangential wake gas experiences centrifugal acceleration:
\begin{equation}
  a_\mathrm{cent} = \frac{u_\mathrm{wake}^2}{R}\,.
  \label{eq:acent}
\end{equation}
The time for gas to traverse the tip gap
$\Delta r = \Delta D / 2$ under this acceleration:
\begin{equation}
  t_\mathrm{cent} = \sqrt{\frac{2\,\Delta r}{a_\mathrm{cent}}}
                   = \sqrt{\frac{2\,\Delta r\, R}{u_\mathrm{wake}^2}}\,.
  \label{eq:tcent}
\end{equation}
For default parameters, $t_\mathrm{cent} \approx \SI{0.04}{s}$, but the
``rotation'' that drives the
centrifugal force is itself decaying. Furthermore, the wake is
\emph{tangential}: it flows \emph{along} the cylindrical wall, not
\emph{into} it. There is no face-on impingement that would deflect
the flow into the axial direction.

\subsection{Axial Coupling Factor ``$\varepsilon$''}

One might posit an azimuthal-to-axial coupling fraction
$\varepsilon$ such that $v_\mathrm{axial} \sim \varepsilon\,
V_\mathrm{mean}$. This parameter has \emph{no basis} in any
identifiable flow mechanism at $\alpha = 0$. All candidate coupling
mechanisms (Ekman pumping, tip clearance, centrifugal wall deflection)
have been shown to be either absent or far too slow.

\subsection{Result}

\begin{equation}
  \boxed{u(z_\mathrm{obs},\, t_\mathrm{obs})
    = 0 \quad \text{(secondary flows)}}
\end{equation}


% ══════════════════════════════════════════════════════════════
\section{Mechanism 6: Vortex Shedding}
\label{sec:vortex}
% ══════════════════════════════════════════════════════════════

When the blade stops, any bound vorticity is shed. For a
\emph{non-lifting} blade ($\alpha = 0$), there is no net circulation,
and hence no starting/stopping vortex of macroscopic strength.

Drag-induced vortex shedding (von K\'arm\'an street) is characterised by
the Strouhal number:
\begin{equation}
  \mathrm{St} = \frac{f\, t_\mathrm{max}}{V}\,,
\end{equation}
where $\mathrm{St} \approx 0.21$ for bluff bodies. But NACA\,0012 at
$\alpha = 0$ is \emph{streamlined}: the wake is thin and coherent
vortex shedding is strongly suppressed. Any residual shed vortices are
confined to the blade plane (within $\sim\delta$ of the surface) and
subject to the same vertical transport limitations established in
Section~\ref{sec:eddies}.

\begin{equation}
  \boxed{u(z_\mathrm{obs},\, t_\mathrm{obs})
    = 0 \quad \text{(vortex shedding)}}
\end{equation}


% ══════════════════════════════════════════════════════════════
\section{Effect of Vessel Confinement}
\label{sec:confinement}
% ══════════════════════════════════════════════════════════════

The preceding analysis treats the gas domain as semi-infinite above the
blade. In reality the propeller operates inside a cylindrical pressure
vessel of height $H$ and inner radius $R_\mathrm{cyl} = R + \Delta r$,
where the radial clearance $\Delta r = \Delta D/2$ is small (of order
\SI{5}{}--\SI{50}{mm}). We now examine whether this confinement
introduces any new mechanism that could transport momentum to
$z_\mathrm{obs}$.

\subsection{Wall Boundary Layers}

The cylindrical wall and the top and bottom endcaps are stationary.
When the blade sweeps, the tangential wake gas flows past these
surfaces and develops wall boundary layers. These boundary layers are
thin: for a characteristic velocity $u_\mathrm{wake} \sim V_\mathrm{tip}/2$
and a flow duration $T_\mathrm{sweep}$, the Stokes penetration depth is
\begin{equation}
  \delta_\mathrm{wall} \sim \sqrt{\nu\, T_\mathrm{sweep}}
    \approx \sqrt{\SI{2.94e-7}{} \times \SI{0.5}{}}
    \approx \SI{0.4}{mm}\,.
\end{equation}
This is confined to the immediate vicinity of the vessel walls and does
not affect the gas at $z_\mathrm{obs}$ above the blade, which is far
from any wall surface.

\subsection{Ekman Spin-Up}

In a confined rotating flow, the Ekman mechanism couples the bulk
rotation to the endcap boundary layers and spins the fluid down over
a characteristic timescale:
\begin{equation}
  t_\mathrm{Ekman} = \frac{H}{\sqrt{\nu\,\Omega}}\,,
  \label{eq:tEkman}
\end{equation}
where $\Omega$ is the representative angular velocity of the gas.
Even using the turbulent viscosity $\nu_t$ (available only during the
first few eddy lifetimes $\sim 2 t_e$), this timescale is
\begin{equation}
  t_\mathrm{Ekman}^{(\mathrm{turb})}
    = \frac{H}{\sqrt{\nu_t\,\Omega}}
    \sim \SI{50}{s}\,,
\end{equation}
far longer than any observation time of interest. The molecular Ekman
time is even longer ($\sim\SI{1000}{s}$). The Ekman pumping velocity
$w_E = \sqrt{\nu\,\Omega}$ is of order \SI{1}{mm/s} (molecular) or
\SI{2}{cm/s} (turbulent), but the turbulent Ekman layer exists for
only $\sim\SI{10}{ms}$, producing a vertical displacement of
$\sim\SI{0.2}{mm}$---negligible compared to $z_\mathrm{obs}$.

\subsection{Tip-Gap Flow (Couette)}

In the radial clearance $\Delta r$ between the blade tip and the
vessel wall, the gas is sheared between the moving blade
(at velocity $V_\mathrm{tip}$) and the stationary wall. The local
Reynolds number based on the gap width is
\begin{equation}
  \mathrm{Re}_\mathrm{gap}
    = \frac{V_\mathrm{tip}\,\Delta r}{\nu}\,.
\end{equation}
For $\Delta r = \SI{5}{mm}$ and $V_\mathrm{tip} \sim \SI{14}{m/s}$,
$\mathrm{Re}_\mathrm{gap} \approx 2.4 \times 10^5$, which is high
enough to be turbulent. However, this Couette-like flow is
\emph{tangential}: the velocity is entirely in the azimuthal
direction, with no mechanism to redirect it vertically. The blade at
$\alpha = 0$ has no pressure difference across its surfaces, so there
is no tip leakage flow. Any vertical component is limited to the
blade-surface boundary layer, which is already accounted for in
Section~\ref{sec:eddies}.

\subsection{Reflected Pressure Waves}

The pressure pulse generated when the blade stops
(Section~\ref{sec:pressure}) reflects off the vessel walls and
endcaps. The longest acoustic path is the vessel diagonal
$\ell_\mathrm{max} = \sqrt{(2R_\mathrm{cyl})^2 + H^2}$, giving a
reverberation time
\begin{equation}
  t_\mathrm{reverb} = \frac{\ell_\mathrm{max}}{c_s}
    \approx \frac{\SI{3.7}{m}}{\SI{178}{m/s}}
    \approx \SI{21}{ms}\,.
\end{equation}
With viscous absorption on every reflection, the acoustic energy is
dissipated within a few tens of milliseconds---well before
$t_\mathrm{obs}$. These waves induce oscillatory velocities (not a
net residual) and leave no persistent flow.

\subsection{Displacement Flow in a Closed Vessel}

In an open domain the displacement (potential) flow around the moving
blade decays as a dipole ($\sim 1/r^2$). In a closed vessel, the
incompressibility constraint modifies the far-field: the gas displaced
by the advancing blade must recirculate within the vessel volume. This
recirculation adjusts at the speed of sound. Once the blade stops, the
source of displacement vanishes, and the potential field readjusts to
zero on the acoustic timescale $R/c_s \approx \SI{9}{ms}$, exactly as
in the unconfined case. The vessel walls provide the return path but
do not sustain any residual flow.

\subsection{Confinement Summary}

No confinement effect introduces a new vertical transport mechanism.
All wall-related phenomena are either tangential (tip-gap Couette),
too slow (Ekman spin-up), too thin (wall BLs), or transient
(reflected pressure waves, displacement readjustment). The six-mechanism
analysis of Sections~\ref{sec:eddies}--\ref{sec:vortex} remains valid
without modification.


% ══════════════════════════════════════════════════════════════
\section{Conclusion}
\label{sec:conclusion}
% ══════════════════════════════════════════════════════════════

Table~\ref{tab:summary} summarises all six mechanisms. Every one yields
zero residual velocity at $z_\mathrm{obs}$ for
$t_\mathrm{obs} \gg t_\mathrm{acoustic}$. The vessel confinement
analysis of Section~\ref{sec:confinement} confirms that no additional
mechanism arises from the finite-volume geometry.

\begin{table}[H]
\centering
\caption{Summary of all mechanisms for residual velocity at
$z_\mathrm{obs}$ above the blade, $t_\mathrm{obs}$ after stopping.
Default parameters: $c = \SI{0.16}{m}$,
$\theta_\mathrm{sweep} = \pi$, $t_\mathrm{rot} = \SI{507}{ms}$ (with drag),
$z_\mathrm{obs} = \SI{5}{mm}$, $t_\mathrm{obs} = \SI{0.5}{s}$.}
\label{tab:summary}
\begin{tabular}{lllc}
\toprule
Mechanism & Characteristic timescale & Reach / Comment
  & $u(z_\mathrm{obs}, t_\mathrm{obs})$ \\
\midrule
Eddy diffusion
  & $2 t_e \approx \SI{5}{ms}$
  & $\delta + \SI{0.7}{mm} = \SI{3.0}{mm}$
  & 0 \\
Bulk swirl
  & $t_\mathrm{visc} \approx \SI{25}{s}$
  & Momentum stays in blade plane
  & 0 \\
Potential flow
  & $t_\mathrm{acoustic} \approx \SI{9}{ms}$
  & Vanishes at speed of sound
  & 0 \\
Pressure pulse
  & $t_\mathrm{decel} \sim \SI{10}{ms}$
  & Transient wave, passes through
  & 0 \\
Secondary flows
  & $t_\mathrm{Ekman} \sim \SI{39}{s}$
  & No lift $\Rightarrow$ no tip clearance
  & 0 \\
Vortex shedding
  & ---
  & Confined to blade plane
  & 0 \\
\bottomrule
\end{tabular}
\end{table}

The gas at $z_\mathrm{obs}$ above the blade surface is \textbf{quiescent}
at time $t_\mathrm{obs}$ after the blade stops, for any
$t_\mathrm{obs}$ much longer than the acoustic clearing time
($\sim\SI{10}{ms}$) and for any $z_\mathrm{obs}$ larger than the
boundary layer thickness ($\sim\SI{2.3}{mm}$ for default parameters).

The fundamental reason is that the NACA\,0012 at zero angle of attack
produces no lift, so all aerodynamic effects are confined to the thin
drag wake within $\delta$ of the blade surface. The dense xenon gas has
an extremely small kinematic viscosity, making all diffusion timescales
very long, while the turbulent eddies are small ($\sim\delta$) and
decay rapidly ($\sim\SI{5}{ms}$) once the blade stops. The cylindrical
vessel walls, endcaps, and radial tip gap introduce no new vertical
transport mechanism (Section~\ref{sec:confinement}): all
confinement-related flows are either tangential, transient, or operate
on timescales orders of magnitude longer than $t_\mathrm{obs}$.

\medskip
\noindent\textbf{Caveat on the gap $z_\mathrm{obs} - \delta$:}\;
The conclusion of zero residual velocity holds provided $z_\mathrm{obs}$
exceeds $\delta$ by at least $\sim\SI{2}{mm}$ (the combined turbulent
and molecular diffusion reach). For parameter choices where
$z_\mathrm{obs} - \delta \lesssim \SI{1}{mm}$ (e.g.\ larger chord
$c = \SI{0.22}{m}$ with $z_\mathrm{obs} = \SI{5}{mm}$), the
exponential tail of the diffusing momentum profile can produce a small
but nonzero residual ($\sim\SI{1}{cm/s}$ at $t_\mathrm{obs} = \SI{1}{s}$).
The numerical PDE solver captures this tail; the companion Python script
should be used to evaluate any specific parameter combination.

\end{document}
